
\documentclass[12pt,a4paper]{article}
\usepackage[a4paper,margin=1.12in]{geometry}
\usepackage[T1]{fontenc}
\usepackage[utf8]{inputenc}
\usepackage{lmodern}
\usepackage{microtype}
\usepackage{amsmath,amssymb,amsthm,mathtools}
\usepackage{enumitem}
\usepackage[hidelinks]{hyperref}
\usepackage{titlesec}
\titleformat{\section}{\large\bfseries}{\thesection}{0.6em}{}
\titleformat{\subsection}{\normalsize\bfseries}{\thesubsection}{0.6em}{}
\newcommand{\webersection}[2]{\section*{\S #1. #2}\addcontentsline{toc}{section}{\S #1. #2}}
\newcommand{\booktitle}[1]{\section*{#1}\addcontentsline{toc}{section}{#1}}
\newcommand{\F}{\mathbb F}
\newcommand{\Fp}{\mathbb F_p}
\newcommand{\Fpp}{\mathbb F_{p^2}}
\newcommand{\Lp}{L_p}
\newcommand{\Ap}{A_p}
\newcommand{\Gp}{G}
\newcommand{\Q}{Q}
\newcommand{\R}{R}
\newcommand{\T}{T}
\newcommand{\eps}{\varepsilon}
\newcommand{\ord}{\operatorname{ord}}
\newcommand{\GL}{\operatorname{GL}}
\newcommand{\SL}{\operatorname{SL}}
\newcommand{\PGL}{\operatorname{PGL}}
\newcommand{\PSL}{\operatorname{PSL}}
\newcommand{\diag}{\operatorname{diag}}
\setlength{\parindent}{0pt}
\setlength{\parskip}{0.76em}
\linespread{1.14}
\title{Heinrich Weber, \emph{Lehrbuch der Algebra}, Volume II\\English Translation -- part 14\\Sections 70--105; Source pp. 271--402}
\author{}
\date{}
\begin{document}
\maketitle
\tableofcontents
\newpage

\webersection{70}{Divisors of the Group \(L_p\) Whose Degree Is Divisible by \(p\)}

It is a problem of the greatest interest to determine all divisors of the group \(L_p\).  On the solution of this problem depends the way in which the group \(L_p\) can be represented by permutations of letters, and therefore the algebraic equations in which these groups can occur.

The cyclic groups contained in \(L_p\) have already been considered in the preceding paragraphs.  We have seen that the degree of a cyclic group is either equal to \(p\), or is a divisor of
\[
 \frac12(p-1)
 \qquad\text{or of}\qquad
 \frac12(p+1),
\]
and every divisor of these numbers also occurs among the degrees of the cyclic groups.  The index of a cyclic divisor can never be smaller than
\[
 \frac12(p^2-1),\qquad p(p+1),\qquad p(p-1),
\]
respectively.

Suppose now that the degree of a divisor \(G\) of \(L_p\) is divisible by \(p\).  Then \(G\) contains a cyclic group of degree \(p\), and by \S 68 we may transform \(G\) so that this cyclic group consists of the substitutions
\[
 T_t:\quad \xi\longmapsto \xi+t,\qquad t=0,1,2,\ldots,p-1 .
\]
If \(G\) contains, besides these, a substitution
\[
 \xi\longmapsto \frac{a\xi+b}{c\xi+d}
\]
in which \(c\ne0\), then by composing it with suitable translations it also contains a substitution interchanging \(0\) and \(\infty\).  From this it follows further that \(G\) contains a substitution of the form
\[
 \xi\longmapsto -\frac{1}{c^2\xi},
\]
and hence, by multiplying with the translations, all substitutions needed to generate the full group \(L_p\).  More explicitly, the substitution of order two obtained in this way, together with the translation \(T_1\), gives a system of generating elements of \(L_p\) as in \S 65.

Consequently, if \(G\) is not the whole group \(L_p\), it can contain only substitutions of the form
\[
 \xi\longmapsto a^2\xi+ab,
\]
or, in matrix notation,
\[
 \begin{pmatrix} a&b\\0&a^{-1}\end{pmatrix},
\]
with the usual identification of proportional matrices.  Conversely, all these substitutions form a group of degree
\[
 \frac12\,p(p-1),
\]
and therefore a divisor of \(L_p\) of index \(p+1\).

Among these groups are contained certain further divisors of greater index.  They are obtained by not allowing \(a\) to run through all admissible values, but only through values of the form \(a^s\), where \(s\) is a divisor of \(\frac12(p-1)\).  The degree of such a group is
\[
 \frac{\frac12\,p(p-1)}{s}.
\]

If the group \(L_p\) is regarded as the group of the linear fractional substitutions
\[
 \eta=\frac{a\xi+b}{c\xi+d}
\]
on the \(p+1\) letters
\[
 \infty,0,1,\ldots,p-1,
\]
then the group just found is the group of the integral linear substitutions
\[
 \eta=a^2\xi+ab .
\]
These are precisely the substitutions leaving the letter \(\infty\) fixed; they therefore give a permutation group on only \(p\) letters.  These are the same special metacyclic groups which were treated in the seventeenth section of the first volume.

The various groups conjugate to \(G\) leave one of the other \(p+1\) letters fixed.  Hence the total number of such groups is \(p+1\).

\webersection{71}{Divisors of the Group \(L_p\) Whose Degree Is Not Divisible by \(p\)}

To find the remaining divisors of \(L_p\), whose degree is not divisible by \(p\), we may proceed step by step along the same path which led us, in the seventh and eighth sections, to the determination of the polyhedral groups.  It is enough here to emphasize the principal points of the derivation, and for the proofs to refer back to that discussion, where the same conclusions were drawn.

For this purpose we again regard \(L_p\), as in \S 65, as a group of linear fractional substitutions, and let the variable \(\xi\) assume all \(p^2+1\) values of the congruence field \(\F_{p^2}\), including \(\infty\).  The advantage thereby gained is that the fixed-point equation
\[
 \xi=\frac{a\xi+b}{c\xi+d}
\]
always has two roots.  These two roots are distinct if we exclude substitutions of \(p\)-th order; such substitutions cannot occur in a group whose degree is not divisible by \(p\), and, by \S 68, they are exactly those for which the two roots coincide.

We call these two roots the \emph{poles} of the substitution.  Let \(G\) be a group of degree \(n\), not divisible by \(p\), which is to be a divisor of \(L_p\).  If in \(G\) the substitutions
\[
 \Theta_1,\Theta_2,\ldots,\Theta_{\nu-1},
\]
and no others, have the same pole \(\alpha\), then we call \(\alpha\) a \(\nu\)-fold pole.

Let \(S\) be any substitution of the group over \(\F_{p^2}\).  If \(L_p\) is transformed by \(S^{-1}\), then \(G\) is transformed into an isomorphic group.  Although the substitutions of the latter group need not have real coefficients, each has two poles, obtained by applying \(S\) to the poles of the original substitution; for if \(\alpha\) is a pole of \(\Theta\), then
\[
 S\Theta S^{-1}S(\alpha)=S\Theta(\alpha)=S(\alpha).
\]
Hence, as in \S 51, one may transform \(G\) so that any chosen substitution has the poles \(0\) and \(\infty\), and so becomes multiplicative.

We now list, in the order used before, the propositions from \S 52 which are needed here.  Additional remarks are added only where the altered hypotheses require them.

If \(\alpha\) is a \(\nu\)-fold pole, then the substitutions
\[
 1,\Theta_1,\Theta_2,\ldots,\Theta_{\nu-1}
\]
form a cyclic group.  Its substitutions may all be transformed into the form
\[
 \Theta(\xi)=\gamma^{\mu\frac{p^2-1}{\nu}}\xi,\qquad
 \mu=0,1,\ldots,\nu-1,
\]
where \(\gamma\) is a primitive element of \(\F_{p^2}\); thus
\(\gamma^{(p^2-1)/\nu}\) is a primitive \(\nu\)-th root of unity in this field.

The two poles of a substitution are of the same multiplicity.  If \(\Theta\) denotes the cyclic group generated by such a substitution, and if \(n=\nu a\), then the poles may be arranged in systems of conjugate poles
\[
 \alpha,\;V_1(\alpha),\;V_2(\alpha),\ldots,V_{a-1}(\alpha),
\]
where the \(V_i\) run through suitable representatives of the group.  Hence all poles of \(G\) split into systems of conjugate poles, and we obtain the same indeterminate equation as in the theory of the polyhedral groups:
\[
 2n-2=a(\nu-1)+a'(\nu'-1)+a''(\nu''-1)+\cdots
      =n\lambda-a-a'-a''-\cdots ,
\]
where \(\lambda\) is the number of systems of conjugate poles.  In \S 52 it was shown that this equation permits only a finite list of solutions.

To obtain the corresponding divisors of the congruence groups, one may put into the formulae for the polyhedral groups, obtained in the eighth section, the roots of unity of the field \(\F_{p^2}\) in place of the complex roots of unity which occur there.  These roots of unity may of course have only degrees dividing \(p^2-1\).  In each case it remains to decide whether the substitutions thus obtained are contained in \(L_p\), or in a group isomorphic with it.

We thus obtain the following cases.  In all of them \(\gamma\) denotes a primitive element of \(\F_{p^2}\).

First, there are the cyclic groups \(C_n\) of degree \(n\):
\[
 C_n:\quad
 \xi\longmapsto \gamma^{\mu\frac{p^2-1}{n}}\xi,\qquad
 \mu=0,1,\ldots,n-1 .
\]
This group is real when \(n\) divides \(\frac12(p-1)\), and only then; for precisely in this case the exponents of \(\gamma\) are divisible by \(p+1\).  The group is contained in \(L_p\), and at the same time in the real form \(A_p\), when the corresponding conjugacy condition is satisfied, that is, precisely when \(n\) divides \(\frac12(p+1)\).  This agrees with \S 68: no cyclic groups other than those whose degree is \(p\), or a divisor of \(\frac12(p-1)\) or of \(\frac12(p+1)\), occur.

In all remaining polyhedral groups one has three systems of exceptional poles.  Thus, secondly, there are the dihedral groups of degree \(2m\).  Using the second representation of the dihedral group from \S 55, the group is written more conveniently for the present purpose.  It is contained in \(L_p\) if
\[
 p\equiv 1\pmod {2m},
\]
and in the imaginary form \(F_p\) if
\[
 p\equiv -1\pmod {2m}.
\]
The distinction is exactly the same as in the cyclic case.

Thirdly, we obtain tetrahedral groups.  The data of the pole multiplicities are
\[
 \nu=2,\qquad \nu'=3,\qquad \nu''=3,
\]
and the degree is \(12\).  In \(\F_{p^2}\) an eighth root of unity always exists.  The formulae of \S 56 therefore give three generating substitutions for the tetrahedral group.  If
\[
 p\equiv 1\pmod 4,
\]
the necessary fourth roots are real, and the group itself is real.  If
\[
 p\equiv 3\pmod 4,
\]
the two conjugate fourth roots lie in the imaginary part of \(\F_{p^2}\); a comparison according as \(p\equiv 3\) or \(7\pmod 8\) shows that the three generators nevertheless belong to the appropriate real or imaginary form.  Thus in every case the group \(L_p\) contains a tetrahedral group.

For example, when \(p=5\), since \(-2\) is a quadratic non-residue modulo \(5\), one may put
\[
 \varepsilon=\sqrt{-2}.
\]
Choosing the primitive roots as in \S 56 gives explicitly a tetrahedral group in \(L_5\).  Written as fractional substitutions, its elements form a divisor of \(L_5\) of index \(5\).

Fourthly, there are octahedral groups.  Since an octahedral group contains a cyclic group of degree \(4\), it can occur only when one of the numbers
\[
 \frac12(p-1),\qquad \frac12(p+1)
\]
is divisible by \(4\), that is,
\[
 p\equiv \pm1\pmod 8 .
\]
Under this condition the formulae of \S 57 give an octahedral group.  This group is real when \(p\equiv 1\pmod8\), and it is imaginary but contained in \(F_p\) when \(p\equiv -1\pmod8\).

The first example is \(p=7\).  In this case \(-1\) is a quadratic non-residue, and one may put \(\varepsilon=\sqrt{-1}=i\).  A convenient primitive element is \(\gamma=2+i\), with conjugate \(\gamma^p=2-i\).  Since \(5\) is a real primitive root modulo \(7\), every non-zero element of \(\F_{7^2}\) is obtained in the form
\[
 \gamma^\mu 5^\nu ,
\]
with the indicated ranges of \(\mu\) and \(\nu\).  The generators produced by the formulae of \S 57 may then be transformed, by the method of \S 69, into real form.  The octahedral group so obtained is a divisor of \(L_7\) of index \(7\).

Finally, the icosahedral group.  For \(p=5\) the group \(L_p\) itself is an icosahedral group.  For other values of \(p\), an icosahedral group can be contained in \(L_p\) only if \(p^2-1\) is divisible by \(5\), i.e.
\[
 p\equiv \pm1\pmod5.
\]
Then the formulae of \S 58 yield the generators of such a group.  If \(p\equiv1\pmod5\), the required fifth roots are real, and the whole group is real.  If \(p\equiv-1\pmod5\), the corresponding roots are conjugate, and the group first obtained lies in the imaginary form \(F_p\); to obtain the icosahedral group in \(L_p\) one has to perform the transformation described in \S 69.

For \(p=11\) one obtains directly a real icosahedral group of index \(11\).  We may choose, for the relevant root, any primitive root of \(11\), for example \(8\); the formulae of \S 58 then give the desired substitutions.

We have thus determined which kinds of divisors can occur in \(L_p\), and we have also shown that all of them actually occur.  We have not pursued here the question how, for each type, all possible divisors are to be obtained.  We add only the following remarks.

If a group \(G\) already found is transformed by a substitution of \(L_p\), one obtains a conjugate divisor.  If, however, one transforms by an imaginary substitution of the group \(F_p\), the transformed group may nevertheless be real and need not be conjugate with \(G\) inside \(L_p\), although both are conjugate as divisors of \(F_p\).  In this way the octahedral and icosahedral groups generally give a second non-conjugate group.  For the tetrahedral group this gives a new group only in the appropriate congruence case.  The examples \(p=5,7,11\) exhibit these alternatives explicitly.

\webersection{72}{Constitution of the Group \(L_7\) of Degree \(168\)}

For later use we must give the group \(L_7\) a particularly explicit form.  In \S 71 it was shown that \(L_7\) contains octahedral divisors of degree \(24\), hence of index \(7\).  It was also shown earlier that \(L_7\) contains divisors of degree \(21\), hence of index \(8\); these are especially important for our purpose.

Among the conjugate divisors of index \(8\), the simplest and most convenient for calculation is the one consisting of substitutions of the form
\[
 \xi\longmapsto a^2\xi+ab .
\]
We arrange the preceding transformation so that this subgroup occurs in precisely this form.  After transforming the generators of the octahedral group found in \S 71, one obtains generators which we shall denote by
\[
 \chi,\quad \omega,\quad \Theta,
\]
of orders \(3,2,4\), respectively.  They generate an octahedral subgroup of \(L_7\).  In Weber's displayed table the twenty-four elements of this subgroup are arranged in six rows, each row containing four elements.  Symbolically they may be written in the form
\[
 \chi^\lambda \omega^\mu \Theta^\nu,\qquad
 \lambda=0,1,2,\quad \mu=0,1,\quad \nu=0,1,2,3,
\]
with the customary relations of the octahedral group.

To obtain the full group \(L_7\) we adjoin an element \(\tau\) of seventh order.  Thus the elements of \(L_7\) are represented in the form
\[
 \tau^q\chi^\lambda\omega^\mu\Theta^\nu,
 \qquad
 q=0,1,\ldots,6,\quad
 \lambda=0,1,2,\quad
 \mu=0,1,\quad
 \nu=0,1,2,3.
\]
No two of these \(168\) expressions represent the same element.  Indeed, no power of \(\tau\) whose exponent is not divisible by \(7\) can lie in the octahedral subgroup, since the octahedral subgroup has no element of seventh order.

The subgroup of degree \(21\) contained in \(L_7\) is represented in this notation by
\[
 \tau^q\chi^\lambda .
\]
In order to describe the composition of the elements it is enough to express, for each \(q\), the products
\[
 \Theta\tau^q,\qquad \omega\tau^q,\qquad \chi\tau^q
\]
again in the normal form just given.  Together with the multiplication table inside the octahedral subgroup this determines the product of any two elements of \(L_7\).

One relation is especially simple:
\[
 \chi\tau=\tau^4\chi.
\]
From it the analogous relations for the powers of \(\chi\) follow.  The remaining products are obtained by straightforward calculation from the octahedral table.  Weber records twelve such relations; they may be reduced to four fundamental relations, from which the others follow with the aid of the octahedral relations and the preceding formula.

The result may be summarized as follows.  Suppose four elements
\[
 \tau,\quad \chi,\quad \omega,\quad \Theta
\]
of orders
\[
 7,\quad 3,\quad 2,\quad 4
\]
are given, and suppose that their compositions satisfy the octahedral relations among \(\chi,\omega,\Theta\), together with the relations obtained by bringing
\(\Theta\tau^q,\omega\tau^q,\chi\tau^q\) into the above normal form.  Then the \(168\) elements
\[
 \tau^q\chi^\lambda\omega^\mu\Theta^\nu
\]
form a simple group of degree \(168\).  The two elements \(\omega\) and \(\tau\) may be taken as generating elements of this group.

Among the divisors of this group one should especially note the octahedral group
\[
 \chi^\lambda\omega^\mu\Theta^\nu
\]
of index \(7\), the group
\[
 \tau^q\chi^\lambda
\]
of degree \(21\) and index \(8\), and the group of degree \(8\) generated by suitable involutions.  The whole group can also be written with the factors in other orders, for example in one of the forms
\[
 \chi^\lambda\omega^\mu\Theta^\nu\tau^q,\qquad
 \omega^\mu\Theta^\nu\chi^\lambda\tau^q,\qquad
 \tau^q\omega^\mu\Theta^\nu\chi^\lambda .
\]

\booktitle{Third Book. Applications of Group Theory}

\webersection{73}{The Resolvents of the Composition Series}

The theorems of general group theory treated in the first book of this volume yield important algebraic consequences when applied to the Galois group of an equation.

Let \(\Omega\) be an arbitrary field taken as the domain of rationality.  Let
\[
 f(x)=0
\]
be an equation in \(\Omega\), without multiple roots, and let \(P\) be its Galois group.  Suppose \(P\) has a normal divisor \(Q\).  Let \(n\) denote the degree of \(P\), and let \(j\) be the index of \(Q\), so that \(n/j\) is the degree of \(Q\).

In the first volume it was shown that the group of \(f(x)\) is reduced from \(P\) to \(Q\) by adjoining a root of an irreducible normal equation of degree \(j\).  This auxiliary equation of degree \(j\) was called the partial resolvent.

The Galois group of this partial resolvent is obtained as follows.  Let \(\psi\) be a function of the roots of \(f(x)\) belonging to the group \(Q\), and decompose \(P\) into right cosets
\[
 P=Q+Qa+Qb+Qc+\cdots ,
\]
where \(a,b,c,\ldots\) are certain permutations in \(P\).  Then all elements of the coset \(Qa\) carry \(\psi\) into one and the same function \(\psi_a\), and the Galois group of the resolvent consists of the substitutions
\[
 (\psi,\psi_a),\quad(\psi,\psi_b),\quad(\psi,\psi_c),\ldots .
\]
If two such substitutions are composed, one must note that
\[
 Qa\cdot Qb=Qab,
\]
and hence
\[
 (\psi,\psi_a)(\psi,\psi_b)=(\psi,\psi_{ab}).
\]
Thus these substitutions compose in exactly the same way as the cosets compose.  It follows that:

\begin{quote}
The Galois group of the partial resolvent belonging to \(Q\) is isomorphic to the quotient group \(P/Q\).
\end{quote}

Now take any composition series of \(P\), with the corresponding series of indices:
\[
 P,\ P_1,\ P_2,\ldots,P_{\mu-1},\ 1,
 \qquad
 j_1,j_2,\ldots,j_\mu .
\]
By what has just been said, the group of the equation \(f=0\) is reduced from \(P\) to \(P_1\) by adjoining a root of a normal equation of degree \(j_1\).  It is then reduced from \(P_1\) to \(P_2\) by adjoining a root of a normal equation of degree \(j_2\), and so on; finally the equation is completely solved by adjoining a root of a normal equation of degree \(j_\mu\).

The theorem on the invariance of the index series therefore receives a new algebraic meaning:

\begin{quote}
To solve the equation \(f=0\), one has to adjoin successively one root of a normal equation of each of the degrees
\[
 j_1,j_2,\ldots,j_\mu .
\]
The order in which these degrees occur may be changed, but the total collection of them cannot be changed.
\end{quote}

The numbers \(j_1,j_2,\ldots,j_\mu\) are therefore connected more deeply with the nature of an equation, or more generally with the fields defined by it, than the degree of the equation itself.  The degree may be altered in many ways by transforming the equation and taking functions of the roots as new unknowns; but the numbers
\[
 j_1,j_2,\ldots,j_\mu
\]
are always retained.  They are to be regarded as genuine invariants of the normal field defined by the equation.  The degree \(n\) of the group \(P\) itself is among these invariants only as their product:
\[
 n=j_1j_2\cdots j_{\mu-1}j_\mu .
\]

In general, in a composition series of \(P\), each term is a normal divisor only of the term immediately preceding it.  It may happen, however, that certain terms are normal divisors of earlier terms as well.  Particularly important are those terms which are normal divisors of \(P\) itself, and consequently of all preceding terms in the composition series.

Let \(P_\nu\) be a term of the series which is also a normal divisor of \(P\).  The index of \(P_\nu\) in \(P\) is
\[
 j_1j_2\cdots j_\nu,
\]
and this product is the degree of the partial resolvent which reduces the group \(P\) to \(P_\nu\).  Denote this resolvent by
\[
 \chi(y)=0.
\]
The group of this resolvent, which is a normal equation, is \(P/P_\nu\).  A composition series for this quotient is obtained from the original one:
\[
 P/P_\nu,\quad P_1/P_\nu,\quad P_2/P_\nu,\ldots,\quad
 P_{\nu-1}/P_\nu,\quad 1,
\]
with the index series
\[
 j_1,j_2,\ldots,j_{\nu-1},j_\nu .
\]

We shall need one further consequence.  If \(Q\) is any normal divisor of \(P\), then, by \S 7, a composition series of \(P\) may be chosen in which \(Q\) occurs.  If \(R\) is another normal divisor of \(P\), and at the same time a divisor of \(Q\), the composition series may be arranged so that both \(Q\) and \(R\) occur:
\[
 P,\ Q',\ Q'',\ldots,Q,\ldots,R .
\]
From this one obtains the quotient series
\[
 P/Q,\quad Q'/Q,\quad Q''/Q,\ldots,1
\]
and
\[
 P/R,\quad Q'/R,\quad Q''/R,\ldots,Q/R,\ldots .
\]
The index of \(Q'/Q\) in \(P/Q\) is the same as the index of \(Q'\) in \(P\), and hence the same as the index of \(Q'/R\) in \(P/R\); the same holds for the following terms.  Therefore:

\begin{quote}
If \(Q\) and \(R\) are normal divisors of \(P\), and \(R\) is a divisor of \(Q\), then the index series of \(P/Q\) is a part of the index series of \(P/R\).
\end{quote}

If the indices \(j_1,j_2,\ldots,j_\nu\) are all prime numbers, the solution of the resolvent
\[
 \chi(y)=0
\]
is reducible to the solution of a chain of cyclic equations of these prime degrees.  Such a resolvent is therefore metacyclic in the sense established in the seventeenth section of the first volume, where metacyclic equations were identified with the equations otherwise called algebraically solvable.

An irreducible equation \(f(x)=0\) is metacyclic if the index series of its group consists entirely of prime numbers.  In the first volume the conditions for metacyclic equations of prime degree were studied.  We now have to carry out this investigation in the general case, where the degree of the equation may be arbitrarily composite.

\webersection{74}{Metacyclic Equations}

Let \(f(x)=0\) be an irreducible equation of degree \(m\), and let \(P\) be its Galois group.  Then \(P\) is transitive.  Let a composition series with index series be
\[
 P,\ P_1,\ P_2,\ldots,P_{\mu-1},\ 1.
\]
Since the last group \(1\) is intransitive, an intransitive group must occur for the first time somewhere in the series.  Let \(P_\lambda\) be this first intransitive group.  Then all following groups
\[
 P_{\lambda+1},P_{\lambda+2},\ldots
\]
are also intransitive, since they are divisors of \(P_\lambda\).

Recall the theorems of the first volume, \S 158.  Since \(P_\lambda\) is a normal divisor of \(P_{\lambda-1}\), those theorems imply that \(P_{\lambda-1}\) is imprimitive.  After the necessary adjunctions have reduced the group of \(f(x)=0\) to \(P_{\lambda-1}\), a further adjunction of a root of a normal equation of degree \(j_\lambda\) causes \(f(x)\) to split into several irreducible factors
\[
 f(x)=f_1(x)f_2(x)\cdots f_h(x),
\]
all of equal degree.  The number \(h\) of these factors is a divisor of \(j_\lambda\).

Let \(m\) be the degree of \(f(x)\), and \(m_1\) the degree of \(f_1(x)\).  Then
\[
 m=hm_1.
\]
The equations
\[
 f_1=0,\ f_2=0,\ldots,f_h=0
\]
all have the same group; denote it by \(Q\).  If
\[
 \alpha,\alpha_1,\ldots,\alpha_{m_1-1}
\]
are the roots of \(f_1=0\), then \(Q\) is obtained by taking the permutations of these \(\alpha\)'s induced by \(P_\lambda\).

Several distinct permutations in \(P_\lambda\) may induce the same permutation of the roots \(\alpha\).  If \(\pi_1,\pi_2\) are two such permutations, then
\[
 \pi_2\pi_1^{-1}
\]
leaves the \(\alpha\)'s fixed.  Conversely, multiplying any one of them by a permutation of \(P_\lambda\) which fixes the \(\alpha\)'s gives another permutation inducing the same element of \(Q\).  Thus every permutation of the \(\alpha\)'s arises the same number of times from \(P_\lambda\).  If this number is \(p_0\), and if \(q\) is the degree of \(Q\), then
\[
 p_\lambda=p_0q.
\]
Thus the degree of \(P_\lambda\) is a multiple of the degree of \(Q\).

But \(Q\) is transitive, since \(f_1(x)\) is irreducible.  Therefore, by the theorem of the first volume that the degree of a transitive permutation group is divisible by the number of letters it permutes, \(q\) is divisible by \(m_1\).  Consequently \(p_\lambda\) is divisible by \(m_1\).  Since \(P_\lambda\) is the first intransitive group, \(P_{\lambda-1}\) is transitive, and its degree is divisible by \(m=hm_1\).  Thus the index
\[
 j_\lambda=\frac{|P_{\lambda-1}|}{|P_\lambda|}
\]
is divisible by \(h\), in agreement with the preceding conclusion that \(h\) divides \(j_\lambda\).

Suppose now that, besides \(j_\lambda\), the number \(m\) has another prime divisor \(p\).  Then one of the subsequent indices
\[
 j_\lambda,j_{\lambda+1},\ldots,j_\mu
\]
must be divisible by \(p\), and hence they cannot all be equal to \(j_\lambda\).  It follows from the theorem of \S 7 that the composition series of \(P_\lambda\) may be chosen in such a way that among
\[
 P_\lambda,P_{\lambda+1},\ldots,P_{\mu-1}
\]
one of the groups, say \(P_\nu\), is a normal divisor of \(P\).  This \(P_\nu\), being a divisor of the intransitive group \(P_\lambda\), is itself intransitive.  Hence \(P\) itself must be imprimitive.

We have proved:

\begin{quote}
If more than one distinct prime divides the degree of an irreducible equation, then this equation can become reducible by successive adjunctions of roots of cyclic equations only if it is imprimitive.
\end{quote}

Assuming such a reduction possible, we can say a little more about its form.  Let \(P_\nu\) be the first group in the series after \(P_\lambda\) which is a normal divisor of \(P\).  Then, with a suitable arrangement of the composition series, the relevant indices
\[
 j_\lambda,j_{\lambda+1},\ldots,j_\nu
\]
are all equal to the same prime.  The resolvent which reduces \(P\) to \(P_\nu\) has therefore a metacyclic group.

Let
\[
 A,B,\ldots,S
\]
be the systems of intransitivity of \(P_\nu\), and let their number be \(s\).  Put
\[
 m=rs.
\]
After adjoining a root of the corresponding resolvent, \(f(x)\) splits into \(s\) factors of degree \(r\).  These systems are systems of imprimitivity of \(P\).

Let \(Q\) be the totality of all permutations in \(P\) which leave the individual systems \(A,B,\ldots,S\) in place, permuting only the elements inside each system.  As was shown in the first volume, \(Q\) is a normal divisor of \(P\).  By adjoining the roots of an auxiliary equation
\[
 \varphi(y)=0
\]
of degree \(s\), the equation \(f(x)\) splits into \(s\) factors of degree \(r\), each associated with one of the systems \(A,B,\ldots,S\):
\[
 f(x)=f(x,y_1)f(x,y_2)\cdots f(x,y_s).
\]
Since this auxiliary equation reduces the group \(P\) to \(Q\), its group is isomorphic to \(P/Q\).  But the intransitive group \(P_\nu\) is a divisor of \(Q\), because the systems are not moved by \(P_\nu\).  Applying the theorem of \S 73 to the three groups \(P,Q,P_\nu\), and using that the index series of \(P/P_\nu\) consists of prime numbers, we see that the same is true of \(P/Q\).  Hence the auxiliary equation \(\varphi(y)=0\) is metacyclic.

Therefore:

\begin{quote}
If an irreducible equation whose degree contains more than one prime factor splits into factors by successive adjunctions of radicals, then a splitting into \(s\) factors of degree \(r\) is brought about by adjoining the roots of a metacyclic equation of degree \(s\).
\end{quote}

As a special case this contains Abel's theorem: an irreducible equation whose degree \(m\) is not a power of a prime can be solvable by radicals only if, by adjoining the roots of a solvable equation of lower degree dividing \(m\), it first decomposes into factors.

This greatly simplifies the question of solving or even reducing equations by radicals.  Further investigation may be restricted to equations whose degree is a prime or a prime power, since all other cases are brought back to these by repeated application of the theorem.

For example, the question of all irreducible equations of degree \(6\) solvable by radicals has an immediate answer: to obtain all metacyclic equations of degree \(6\) over a field \(\Omega\), adjoin a square root to \(\Omega\) and form all cubic equations in the enlarged field; or adjoin the root of a cubic equation and form all quadratic equations in the enlarged field.

As an example one may cite the equation treated by Hesse, on which the circular sections of a quadric surface not referred to its principal axes depend.  This problem is solved by square roots after the principal axes have first been determined by the solution of a cubic equation.

\webersection{75}{Metacyclic Equations Whose Degree Is a Power of a Prime}

After the preceding theorems, the interest of the further investigation of metacyclic equations is concentrated chiefly on the case in which the degree of the equation is a power
\[
 p^k
\]
of a prime number \(p\).  We suppose the equation irreducible, and restrict ourselves to primitive equations.  The solution of an imprimitive equation reduces to the successive solution of two or more primitive equations whose degrees are again powers of \(p\), and which, if the original equation is metacyclic, must themselves be metacyclic.

Let \(P\) be the group of an irreducible primitive metacyclic equation
\[
 f(x)=0
\]
of degree \(p^k\).  By the theorem of the first volume, not only \(P\) itself, but every normal divisor of \(P\) different from the identity group, must be transitive.  We now apply the theorem of \S 8: \(P\) possesses a normal divisor \(Q\) all of whose elements commute.  If several such divisors exist, choose \(Q\) of the smallest possible degree greater than \(1\).

This group \(Q\) is still transitive.  It cannot have a proper divisor, different from the identity, which is at the same time a normal divisor of \(P\); otherwise that divisor could have been chosen in place of \(Q\).

After the appropriate adjunctions reduce the group of our equation from \(P\) to \(Q\), the equation \(f(x)=0\) has become an abelian equation over the enlarged rationality domain.  Since it has remained irreducible, the theorem of the first volume implies that the degree of the equation is equal to the degree of its group.  Thus \(Q\) is an abelian group of degree
\[
 p^k .
\]
The orders of all elements of \(Q\) are therefore powers of \(p\).

We now show that, apart from the identity, only elements of order \(p\) can occur in \(Q\).  Suppose that \(p^\lambda\) is the highest order occurring among the elements of \(Q\), with \(\lambda>1\).  Then all elements of \(Q\) whose orders divide \(p^{\lambda-1}\) form a divisor \(Q'\) of \(Q\), distinct both from \(Q\) and from the identity group.  This divisor must be a normal divisor of \(P\).  Indeed, if \(\pi\in Q'\) and \(\rho\in P\), then
\[
 \rho^{-1}\pi\rho
\]
has the same order as \(\pi\), and since \(Q\) is normal in \(P\), it again belongs to \(Q\); hence it belongs to \(Q'\).  This contradicts the minimal choice of \(Q\).  Therefore \(\lambda=1\).

Thus \(Q\) is an elementary abelian group: every non-identity element has order \(p\).

\webersection{76}{Representation of the Abelian Group \(Q\)}

The preceding paragraph has shown that in a metacyclic group \(P\), when \(P\) occurs as the Galois group of a primitive irreducible equation \(f(x)=0\) of degree
\[
 n=p^k,
\]
there must be contained, as a normal divisor, an abelian group \(Q\) of degree \(p^k\), and that \(Q\) has no elements except the identity and elements of order \(p\).

By \S 9 this group \(Q\) may be represented by a basis consisting of \(k\) elements
\[
 A_1,A_2,\ldots,A_k
\]
of order \(p\), so that every element of \(Q\) is uniquely of the form
\[
 A_1^{z_1}A_2^{z_2}\cdots A_k^{z_k},
\]
where \(z_1,z_2,\ldots,z_k\) independently run through complete systems of residues modulo \(p\).

We now pass to the corresponding permutation group.  Choose an arbitrary root \(x\) of \(f(x)\), and denote by
\[
 [z_1,z_2,\ldots,z_k]
\]
the root into which \(x\) is carried by the permutation
\[
 A_1^{z_1}A_2^{z_2}\cdots A_k^{z_k}.
\]
This notation is to remain unchanged if the \(z_i\) are altered by multiples of \(p\).  The root from which we started is then
\[
 [0,0,\ldots,0],
\]
and every root of \(f(x)\) is represented once and only once by such a symbol.

If \(x\) is carried by a permutation \(A'\) into \(x'\), and \(x'\) by \(A''\) into \(x''\), then \(x\) is carried by \(A'A''\) into \(x''\).  It follows that the permutation
\[
 A=A_1^{a_1}A_2^{a_2}\cdots A_k^{a_k}
\]
carries
\[
 [z_1,z_2,\ldots,z_k]
\]
into
\[
 [z_1+a_1,\ z_2+a_2,\ldots,z_k+a_k],
\]
all congruences being taken modulo \(p\).  In this way the whole group \(Q\) is obtained by allowing
\[
 a_1,a_2,\ldots,a_k
\]
to run independently through complete residue systems modulo \(p\).

The group \(Q\), whose construction has now been fixed, must be a normal divisor of \(P\).  From this normality an explicit general form can be derived for all permutations of \(P\).  For this we need first a lemma on the analytic representation of permutations.

\webersection{77}{Analytic Representation of the Permutations}

\textbf{Lemma.}  Suppose an arbitrary permutation of the symbols
\[
 [z_1,z_2,\ldots,z_k]
\]
carries \(z_1,z_2,\ldots,z_k\) into
\[
 z_1',z_2',\ldots,z_k'.
\]
Then one may always write
\[
 z_i'=\varphi_i(z_1,z_2,\ldots,z_k)\pmod p
 \qquad (i=1,\ldots,k),
\]
where each \(\varphi_i\) is an integral rational polynomial in the variables
\[
 z_1,z_2,\ldots,z_k
\]
with integral coefficients, and where no variable occurs to a degree greater than \(p-1\).

This is a generalization of the corresponding theorem in \S 180 of the first volume, where the case \(k=1\) was proved.  The general proof rests on complete induction.  We first note the following.

In the formulae above the arguments \(z_i\) have to run through all combinations of \(k\) residue systems modulo \(p\), hence through \(p^k\) systems.  Each of the functions \(\varphi\) has \(p^k\) undetermined coefficients.  If the values \(z_i'\) are prescribed, we therefore obtain \(kp^k\) linear congruences for the same number of unknown coefficients.  It remains to show that these congruences are independent.

It is enough to consider each function \(\varphi\) separately.  Suppose
\[
 z'=\varphi(z_1,z_2,\ldots,z_k)\pmod p .
\]
For each of the \(p^k\) arguments the corresponding value of \(z'\) is known.  We have to determine the coefficients of \(\varphi\), assuming that the assertion has already been proved for functions of fewer variables.

Arrange \(\varphi\) according to the powers of \(z_1\):
\[
 z'=\chi_0 z_1^{p-1}+\chi_1 z_1^{p-2}+\cdots+\chi_{p-1}\pmod p,
\]
where the coefficients \(\chi_0,\chi_1,\ldots,\chi_{p-1}\) are themselves polynomials of the same kind in
\[
 z_2,\ldots,z_k.
\]
Now keep one combination of \(z_2,\ldots,z_k\) fixed and let \(z_1\) run through \(0,1,\ldots,p-1\).  We then obtain a system of \(p\) linear congruences whose determinant is not divisible by \(p\).  As in the one-variable case of the first volume, this system determines the values of
\[
 \chi_0,\chi_1,\ldots,\chi_{p-1}
\]
for the fixed combination of \(z_2,\ldots,z_k\).  Hence, for every combination of the remaining variables, the values of each \(\chi_i\) are known.  By the induction hypothesis the coefficients of the functions \(\chi_i\) may therefore be determined.  The lemma is proved.

Consequently every permutation of the roots may be written
\[
 [z_1,\ldots,z_k]\longmapsto
 [\varphi_1(z_1,\ldots,z_k),\ldots,\varphi_k(z_1,\ldots,z_k)].
\]
If the new variables \(z_i'\) are themselves written as polynomial functions of the old variables, the same symbol may be rewritten with the \(z_i'\) in place of the \(z_i\).  This gives the rule for composing permutations: substitution of one system of polynomial functions into the other.

\webersection{78}{Representation of the Metacyclic Group \(P\)}

The group \(Q\) is to be a normal divisor of \(P\).  Thus, if \(A\) is a permutation in \(Q\) and \(B\) a permutation in \(P\), a second permutation \(A'\) in \(Q\) must be determined so that
\[
 AB=BA'.
\]
Using the abbreviated notation of the preceding paragraph, write
\[
 A:z\mapsto z+a,\qquad
 A':z\mapsto z+a',\qquad
 B:z\mapsto \varphi(z).
\]
Then
\[
 AB:z\mapsto \varphi(z+a),\qquad
 BA':z\mapsto \varphi(z)+a',
\]
and the condition becomes
\[
 \varphi(z+a)\equiv \varphi(z)+a'\pmod p.
\]
Written without abbreviation this is the system
\[
\begin{aligned}
 \varphi_1(z_1+a_1,z_2+a_2,\ldots)
     &\equiv \varphi_1(z_1,z_2,\ldots)+a_1',\\
 \varphi_2(z_1+a_1,z_2+a_2,\ldots)
     &\equiv \varphi_2(z_1,z_2,\ldots)+a_2',\\
 &\ \vdots
\end{aligned}
\]
for all choices of \(z\) and \(a\).  These congruences show that the differences of the functions \(\varphi_i\) are constant under all translations.  It follows that each \(\varphi_i\) is a linear function plus a constant:
\[
 z_i'=a_{i1}z_1+a_{i2}z_2+\cdots+a_{ik}z_k+b_i\pmod p.
\]
The determinant \((a_{ij})\) cannot vanish, since \(B\) is a permutation of all \(p^k\) symbols.

Thus every group \(P\) of the kind under consideration is contained in the full linear congruence group
\[
 R:\qquad
 z'=Az+b
\]
over the field of residues modulo \(p\).  Conversely, this affine group contains the translation group \(Q\) as a normal divisor.

Let \(S\) denote the group of homogeneous linear substitutions
\[
 z_i'=a_{i1}z_1+a_{i2}z_2+\cdots+a_{ik}z_k,\qquad
 \det(a_{ij})\ne0,
\]
taken modulo \(p\).  The group \(R\) is then the product
\[
 R=SQ,
\]
and every element of \(R\) may be written in one and only one of the two forms
\[
 \pi=\delta\gamma=\gamma'\delta,
\]
where \(\delta\in S\) and \(\gamma,\gamma'\in Q\).  The translation \(\gamma'\) is obtained from \(\gamma\) by the homogeneous substitution \(\delta\).

The group \(Q\) is a normal divisor of \(R\).  If \(R_1\) is a divisor of \(R\) containing \(Q\), then
\[
 R_1=S_1Q
\]
for a divisor \(S_1\) of \(S\), and the degree of \(R_1\) is the product of the degrees of \(S_1\) and \(Q\).  If \(R_2\) is a normal divisor of \(R_1\), and if \(R_2\) also contains \(Q\), then
\[
 R_2=S_2Q
\]
where \(S_2\) is a normal divisor of \(S_1\).  Conversely, if \(S_2\) is a normal divisor of \(S_1\), then \(S_2Q\) is a normal divisor of \(S_1Q\).

We therefore obtain:

\begin{quote}
If \(P\) is the Galois group of a primitive irreducible metacyclic equation of degree \(p^k\), then \(P\) is representable in the form
\[
 P=TQ,
\]
where \(T\) is a metacyclic group contained in the homogeneous congruence group \(S=\GL(k,p)\).
\end{quote}

The problem of finding all these metacyclic equations is thus reduced to that of finding all metacyclic divisors of the homogeneous congruence group \(S\).

As an example take \(p=3,\ k=2\), and ask for metacyclic equations of degree \(9\).  The group \(S\) is the group of linear substitutions
\[
 \begin{pmatrix}a&b\\ c&d\end{pmatrix}
\]
modulo \(3\).  In \S 66 the corresponding projective group was considered.  There we identified opposite matrices
\[
 \begin{pmatrix}a&b\\ c&d\end{pmatrix},
 \qquad
 \begin{pmatrix}-a&-b\\ -c&-d\end{pmatrix},
\]
and also restricted first to determinant \(1\).  In the present group \(S=\GL(2,3)\) the two matrices are distinct, and the matrices of determinant \(-1\) must be included.  Thus the degree of the group is \(48\).

The group denoted previously by \(G\) is enlarged by changing all signs into a group of degree \(8\); the group denoted by \(E\) is enlarged into a group of degree \(24\), in which the enlarged \(G\) is a normal divisor.  The whole group \(S\) is metacyclic.  A composition and index series is
\[
 S,\ S_1,\ S_2,\ S_3,\ S_4,
 \qquad
 2,\ 3,\ 2,\ 2.
\]
Accordingly, all equations of degree \(9\) whose group is the full linear congruence group are metacyclic.

\webersection{79}{The Ternary Linear Congruence Group for the Modulus \(2\)}

Many interesting relations appear if one considers the ternary linear congruence group \(R\) for the modulus \(2\).  It may be regarded as a transitive permutation group on eight elements and must contain the groups of metacyclic equations of degree \(8\).  According to what has been proved above, the first object is the homogeneous congruence group \(S\), consisting of the matrices
\[
 \delta=
 \begin{pmatrix}
 a&b&c\\
 a_1&b_1&c_1\\
 a_2&b_2&c_2
 \end{pmatrix}
\]
whose entries are \(0\) or \(1\) modulo \(2\), and whose determinant is non-zero, hence equal to \(1\) modulo \(2\).

The degree of \(S\) is obtained as follows.  The first row can have any of the seven non-zero forms
\[
 (1,0,0),\ (0,1,0),\ (0,0,1),\ (0,1,1),\ (1,0,1),\ (1,1,0),\ (1,1,1).
\]
For the first choice \((1,0,0)\), the determinant is reduced to
\[
 b_1c_2-c_1b_2,
\]
and this can be made equal to \(1\) in six ways.  The remaining entries in the first column can then be chosen in four ways.  Thus there are \(24\) matrices with first row \((1,0,0)\), and the same number for each first row.  Therefore
\[
 |S|=7\cdot24=168.
\]

The number \(168\) has already occurred as the degree of the group \(L_7\).  This suggests that \(S\) and \(L_7\) are isomorphic.  If this conjecture is correct, the investigation of \(S\) is greatly simplified, because the divisors of \(L_7\) are already known.

To verify it, we choose generating elements \(\tau,\omega,\Theta,\chi\) of \(S\) so that they satisfy the characteristic relations of Theorem I in \S 72.  This can be done in several ways.  A systematic procedure is lacking, but the choice is readily found by trial.

One first chooses an element \(\tau\) of seventh order and forms its period.  A convenient matrix is
\[
 \tau=
 \begin{pmatrix}
 1&1&1\\
 1&0&1\\
 1&0&0
 \end{pmatrix},
\]
for which the successive powers give the desired cycle of length \(7\).  Next one selects, among the \(21\) elements of second order, an element suitable for \(\omega\), for instance
\[
 \omega=
 \begin{pmatrix}
 1&0&0\\
 0&1&0\\
 0&1&1
 \end{pmatrix}.
\]
The formulae of \S 72 then determine the elements \(\Theta\) and \(\chi\).  A direct calculation verifies the relations of Theorem I in \S 72.  Hence \(S\) is isomorphic to \(L_7\).

It follows that the group \(S=\GL(3,2)\), like \(L_7\), is simple.

\webersection{80}{Resolvents of the Equation of Eighth Degree}

We now apply this to equations of degree \(8\).  Let \(Q\) again denote the elementary abelian group of degree \(8\), consisting of the translations
\[
 (z_1,z_2,z_3)\longmapsto(z_1+a_1,z_2+a_2,z_3+a_3)
 \qquad(\bmod 2),
\]
and let \(S\) be the ternary congruence group of \S 79.  Put
\[
 R=SQ,
\]
so that \(R\) is the full ternary affine congruence group modulo \(2\).

Consider a system of eight independent quantities
\[
 X_{z_1,z_2,z_3},
\]
where the indices \(z_1,z_2,z_3\) are taken modulo \(2\).  From them form the seven Lagrange resolvents
\[
 \Psi_{\xi_1,\xi_2,\xi_3}
   =\sum_{z_1,z_2,z_3}
     (-1)^{\xi_1z_1+\xi_2z_2+\xi_3z_3}
       X_{z_1,z_2,z_3},
\]
where \((\xi_1,\xi_2,\xi_3)\) runs through the seven non-zero triples modulo \(2\).

If a translation in \(Q\) is applied to the indices of the \(X\)'s, then
\[
 X_z\longmapsto X_{z+a}.
\]
Consequently
\[
 \Psi_\xi\longmapsto
 (-1)^{\xi\cdot a}\Psi_\xi .
\]
Thus the functions \(\Psi\) are changed by the permutations of \(Q\) only in sign, and their squares remain unchanged under \(Q\).

It remains to determine the effect of a homogeneous substitution \(\delta\in S\) on the functions \(\Psi\).  Let \(q\) be the transpose of \(\delta\).  If
\[
 z'=\delta z,
\]
then the scalar product transforms by
\[
 \xi\cdot z'=(q\xi)\cdot z.
\]
Thus applying the permutation \(\delta\) to the indices of the \(X\)'s is equivalent, in the functions \(\Psi\), to applying \(q^{-1}\) to the indices \(\xi\).  Therefore:

\begin{quote}
The application of a permutation \(\delta\in S\) to the quantities \(X\) induces on the seven resolvents \(\Psi_\xi\) the inverse transposed substitution on the index triples \(\xi\).
\end{quote}

\webersection{81}{The Triple Systems of the Resolvents}

We now write \(\Psi_\xi\) simply as \(\Psi\), when the index is clear.  Besides the squares of the seven functions, certain products of them are also invariant under the translation group \(Q\).  The first is the product of all seven:
\[
 A=\prod_{\xi\ne0}\Psi_\xi .
\]

There are also products of three functions
\[
 \Psi_\xi\Psi_\eta\Psi_\zeta
\]
and their reciprocal companions, provided the indices satisfy
\[
 \xi+\eta+\zeta=0\pmod2 .
\]
We shall call such a system of three functions a \emph{triple}; likewise, the three non-zero index systems \(\xi,\eta,\zeta\) satisfying this relation will be called a triple.

Altogether there are exactly seven such triples.  No two of the three systems in a triple can be equal, and any two determine the third.  Thus, if the first two systems are chosen among the seven possibilities, every triple is counted six times; the number is therefore \(7\cdot6/6=7\).

To characterize the triples, introduce a non-zero triple
\[
 \alpha=(\alpha_1,\alpha_2,\alpha_3)
\]
of residues modulo \(2\).  The congruence
\[
 \alpha_1\xi_1+\alpha_2\xi_2+\alpha_3\xi_3\equiv0\pmod2
\]
is satisfied by exactly three non-zero systems \(\xi\).  These three form a triple.  Thus the triple belonging to \(\alpha\) is characterized by
\[
 \alpha\cdot\xi=0,\qquad
 \alpha\cdot\eta=0,\qquad
 \alpha\cdot\zeta=0.
\]
It follows that the function attached to the triple may be denoted unambiguously by \(v_\alpha\).

As functions of the eight quantities \(X\), the \(v_\alpha\) admit all permutations of \(Q\).  If a substitution \(\delta\in S\) is applied to the indices of the \(X\)'s, then the indices of the \(\Psi\)'s undergo the inverse transposed substitution.  From the defining equations of the triple it follows that the index \(\alpha\) itself undergoes the direct substitution \(\delta\).  Hence:

\begin{quote}
The application of a substitution \(\delta\in S\) to the indices of \(X\) has the same substitution as its effect on the indices \(\alpha\) of the functions \(v_\alpha\).
\end{quote}

Now take a fixed non-zero system \(\alpha\).  The three \(\xi\)'s satisfying
\[
 \alpha\cdot\xi=0
\]
form a triple; the four remaining non-zero systems, together with the complementary parity condition, form a quadrupel.  Each triple determines a quadrupel and conversely; hence there are seven quadrupels as well.

If \(\beta\ne\alpha\), then among the four sums
\[
 \beta\cdot\xi
\]
as \(\xi\) runs through the quadrupel belonging to \(\alpha\), two are even and two are odd.  This is immediate by comparing with the three elements of the complementary triple.  This counting also yields Weber's product relation among \(A\), the four functions \(v\) belonging to a quadrupel, and the corresponding resolvent factor.  In the notation where the quadrupel belonging to \(h\) is
\[
 \alpha,\beta,\gamma,\delta,
\]
one obtains
\[
 \Psi_h^{\,4}=A^2v_\alpha v_\beta v_\gamma v_\delta .
\]

For practical use the seven triples and the seven complementary quadrupels are most conveniently tabulated.  Number the seven non-zero triples
\[
 100,\ 010,\ 001,\ 011,\ 101,\ 110,\ 111
\]
by \(1,2,3,4,5,6,7\).  Then each number determines the three numbers lying in its associated triple, and the complementary four numbers are the associated quadrupel.  This is the combinatorial table used in the following application.

\webersection{82}{Application to Equations of Eighth Degree}

We now put, in the formulae just derived, the roots of an equation of eighth degree for the eight variables \(X_z\).  Suppose its Galois group is
\[
 P=TQ
\]
over the fixed rationality domain, and that this group is contained in the affine congruence group
\[
 R=SQ,
\]
so that \(T\) is a divisor of \(S\).  All functions which admit the permutations of \(P\) are rational.

First among these is the sum of the roots:
\[
 B=\sum_z X_z.
\]
Then there is the product \(A\) defined in \S 81.  Further, the symmetric functions of the seven squares
\[
 \Psi_\xi^2
\]
are rational.  If, for simplicity, we assume for the moment that none of the seven quantities \(\Psi_\xi\) vanishes, then the functions \(v_\alpha\) are also defined.  Not merely the symmetric functions of the \(v_\alpha\)'s, but every function of them which admits the permutations of the group \(T\), is rational.

The seven quantities \(v_\alpha\) are therefore the roots of a rational equation of degree \(7\), with group \(T\).  By using the equation for the sum of the roots and the seven equations expressing the \(\Psi\)'s, one obtains, after applying the product relations of the preceding paragraph, an explicit expression for a root such as \(X_{0,0,0}\).  In Weber's notation it has the form
\[
 8X_{0,0,0}
  =B+\sqrt{v_\alpha v_\beta v_\gamma v_\delta}
      +\sqrt{v_{\alpha'}v_{\beta'}v_{\gamma'}v_{\delta'}}
      +\cdots ,
\]
where the displayed square roots run through the quadrupels appearing in the table of \S 81.

There are fifteen sign changes of the seven square roots under which this expression remains unchanged: the simultaneous change of all seven signs, and the changes of the signs belonging to a triple or to a quadrupel.  Therefore the expression assumes only eight distinct values, however the seven square roots occurring in it are chosen.  These eight values are precisely the eight roots.

\webersection{83}{Metacyclic Equations of Eighth Degree}

We can now determine the primitive metacyclic equations of eighth degree.  First one has to find the metacyclic divisors of \(S=\GL(3,2)\).  The group \(S\) itself, as already observed, is not metacyclic: it is the simple group of degree \(168\).  Its proper divisors were considered in \S\S 70--72, and all proper divisors of \(S\) are metacyclic.

Among them are divisors of degrees \(21\) and \(7\), which we denote by
\[
 T_{21},\qquad T_7,
\]
and also divisors of degree \(24\), together with other divisors whose degrees divide \(24\), which we collect under the symbol \(T_{24}\).

The groups \(T_{24}\) cannot connect the seven non-zero systems
\[
 \xi=(\xi_1,\xi_2,\xi_3)
\]
transitively, because the degree of a transitive permutation group on seven elements must be divisible by \(7\).  We shall show that the permutation groups corresponding to \(T_{24}Q\) are imprimitive.

Adjoin to the seven systems \(\xi\) the eighth system
\[
 (0,0,0),
\]
and denote the eight quantities \(X_z\) by
\[
 0,1,2,3,4,5,6,7.
\]
There are two cases.

In the first case the substitutions of \(T_{24}\) leave two of these eight elements, say \(0,1\), fixed.  In \(Q\) there is exactly one substitution \(\gamma_0\) which interchanges \(0\) and \(1\).  The two substitutions \(1,\gamma_0\) form a group \(Q_0\).  If \(\gamma_1\in Q\) is not in \(Q_0\), then it carries \(0,1\) into two different elements, say \(2,3\).  Similarly \(Q\) is decomposed into four cosets
\[
 Q=Q_0+Q_0\gamma_1+Q_0\gamma_2+Q_0\gamma_3,
\]
and the eight elements are split into four pairs
\[
 0,1;\qquad 2,3;\qquad 4,5;\qquad 6,7 .
\]
Each coset carries the pair \(0,1\) into one of these four pairs.

Now let \(h\) be any substitution of the group \(T_{24}Q\).  Suitable elements of \(T_{24}\) and of \(Q\) may be chosen so that, on multiplying by the representatives above, each of the four pairs is carried into one of the same four pairs.  Thus the group \(T_{24}Q\) is imprimitive, with four systems of imprimitivity.  The two roots \(0,1\) satisfy a quadratic equation whose coefficients depend on the roots of a biquadratic equation.

In the second case the elements of a triple
\[
 1,2,3
\]
and the elements of its complementary quadrupel
\[
 4,5,6,7
\]
are each connected transitively by \(T_{24}\).  Then the eight quantities split into the two quadrupels
\[
 0,1,2,3;\qquad 4,5,6,7 .
\]
The translations in \(Q\) either permute the elements inside one of these quadrupels, or carry the whole quadrupel into the other.

Indeed, write \(1,2,3\) as the three systems
\[
 \xi,\eta,\zeta.
\]
Since they form a triple, there is a non-zero system \(\alpha\) such that
\[
 \alpha\cdot\xi=\alpha\cdot\eta=\alpha\cdot\zeta=0\pmod2.
\]
After adding a translation \(\mu\) to the indices, the three translated systems form the same triple if \(\mu\) belongs to it; otherwise, together with \(\mu\), they form the complementary quadrupel.  Hence any element of \(T_{24}Q\) carries one of the two quadrupels into one of the two quadrupels.  The group is again imprimitive.  After adjoining a square root to the rationality domain, the four quantities
\[
 0,1,2,3
\]
are the roots of a biquadratic equation.

These considerations also exclude, for primitive equations of degree \(8\), the possibility that one of the seven quantities \(\Psi_\xi\) vanish.  The functions \(\Psi_\xi\) are rational functions of the roots \(X_z\).  If one of them vanished, then the equation \(\Psi_\xi=0\) would be preserved by every permutation of the group of the equation.  By the theorem of \S 80, either all seven \(\Psi_\xi\) would have to vanish, in which case the roots \(X_z\) themselves would be rational, or the induced group on the seven non-zero index systems would be intransitive.  Then the degree of \(T\) would not be divisible by \(7\), and the corresponding group would fall under the imprimitive cases just discussed.

Thus in the primitive metacyclic case one is led to the divisors whose action on the seven non-zero indices has degree divisible by \(7\), namely the groups \(T_7\) and \(T_{21}\).  The preceding resolvent construction gives the associated equations of degree \(8\): a rational equation of seventh degree in the quantities \(v_\alpha\), followed by the extraction of square roots in the formula for the eight \(X_z\)'s.

\webersection{84}{Biquadratic Equations}

The case of biquadratic equations is the same construction in two variables and is sufficiently simple to display explicitly.  We have four roots
\[
 X_{0,0},\quad X_{1,0},\quad X_{0,1},\quad X_{1,1},
\]
and the translation group \(Q\) of order \(4\) acts by adding a pair of residues modulo \(2\) to the indices.  The homogeneous group \(S=\GL(2,2)\) has degree \(6\), and acts by the corresponding linear substitutions on the non-zero index pairs.

The three Lagrange resolvents are
\[
\begin{aligned}
 \Psi_{1,0}&=X_{0,0}-X_{1,0}+X_{0,1}-X_{1,1},\\
 \Psi_{0,1}&=X_{0,0}+X_{1,0}-X_{0,1}-X_{1,1},\\
 \Psi_{1,1}&=X_{0,0}-X_{1,0}-X_{0,1}+X_{1,1}.
\end{aligned}
\]
Their squares are invariant under \(Q\), and the group \(S\) permutes them.

Let the three quantities formed from the products of these resolvents be denoted
\[
 v_{1,0},\qquad v_{0,1},\qquad v_{1,1}.
\]
They remain unchanged under \(Q\), and under two generating substitutions of \(S\) they are permuted cyclically and by transposition.  Consequently the elementary symmetric functions of
\[
 v_{1,0},\ v_{0,1},\ v_{1,1}
\]
are rational.  These three quantities are therefore the roots of a cubic equation.  This is the cubic resolvent of the general biquadratic equation.

If we put
\[
 a=X_{0,0}+X_{1,0}+X_{0,1}+X_{1,1},
\]
then the roots are recovered from the formula
\[
 4X_{0,0}
 =
 a+\sqrt{v_{1,0}}\sqrt{v_{0,1}}
  +\sqrt{v_{1,0}}\sqrt{v_{1,1}}
  +\sqrt{v_{0,1}}\sqrt{v_{1,1}}.
\]
By changing the signs of the three square roots in all possible ways this expression gives only four distinct values, and these four values are the four roots of the biquadratic equation.  Conversely, if \(v_{1,0},v_{0,1},v_{1,1}\) are the roots of any cubic equation, the four quantities obtained from this formula satisfy a biquadratic equation with rational coefficients.



\booktitle{Eleventh Section. The Inflection Points of a Curve of the Third Order}

\webersection{85}{Ternary Forms and Algebraic Curves}

In order, in the further course of this exposition, to exhibit some of the many and interesting applications of algebra to geometrical problems without referring the reader to other aids, we shall derive briefly the geometrical foundations required for these applications.  We restrict ourselves to plane geometry, and even there leave aside everything not directly needed for our purpose.

We shall use homogeneous triangular coordinates.  The position of a point \(x\) is first described by its three distances from the sides of a fixed coordinate triangle, these distances being measured in arbitrary units, or, what is the same thing, multiplied by three arbitrary constant factors.  These three quantities are called the coordinates of \(x\), and are denoted by
\[
x_1,\quad x_2,\quad x_3.
\]
In the elementary form of the definition they are subject to one non-homogeneous linear relation, obtained by equating the area of the coordinate triangle \((1,2,3)\) with the sum of the areas of the triangles \((x,2,3)\), \((x,3,1)\), and \((x,1,2)\).  A coordinate changes sign whenever the point passes through the corresponding side of the triangle.

By means of this relation every non-homogeneous expression can be made homogeneous.  We shall therefore enlarge the meaning of the coordinates and regard \(x_1,x_2,x_3\) as independent variables, with the convention that only homogeneous equations are used and that the point is determined not by the three numbers themselves, but by their ratios
\[
x_1:x_2:x_3.
\]
The triple \(0,0,0\) is excluded, and multiplying all three coordinates by the same non-zero factor gives the same point.

A linear substitution
\[
x_i=\sum_{j=1}^3 a_{ij}y_j\qquad (i=1,2,3),
\]
whose determinant
\[
r=\det(a_{ij})
\]
is different from zero, may be interpreted in two ways.  It may be regarded as a projective transformation carrying the point \(y\) to the point \(x\), or as a change of coordinates for one and the same point.  For the present we hold to the latter interpretation.

A ternary form of degree \(n\),
\[
f(x_1,x_2,x_3),
\]
when equated to zero, represents a curve of order \(n\).  We shall call it simply the curve \(f\).  Under the substitution above the form goes over into a form of the same degree,
\[
f(x_1,x_2,x_3)=\varphi(y_1,y_2,y_3),
\]
and the equation \(\varphi=0\) represents the same curve in the new coordinates.  Differentiating this identity gives the transformation formulae for the first partial derivatives:
\[
\frac{\partial\varphi}{\partial y_j}
  =\sum_i a_{ij}\frac{\partial f}{\partial x_i}.
\]
Thus properties expressed by the simultaneous vanishing of such derivatives are invariant under projective change of coordinates.

\webersection{86}{Singular Points. Inflection Points. Bitangents}

If at a point \(x\) the three equations
\[
f_{x_1}=0,\qquad f_{x_2}=0,\qquad f_{x_3}=0
\]
are simultaneously satisfied, then by Euler's theorem for homogeneous functions,
\[
nf=x_1f_{x_1}+x_2f_{x_2}+x_3f_{x_3},
\]
the point lies on the curve \(f=0\).  The equations just written are transformed into equations of the same kind under every linear change of coordinates.  A point so characterized is therefore defined by an invariant condition.  Such points are called singular points of the curve.

Not every curve of order \(n\) has singular points.  In general their existence imposes a condition on the coefficients of \(f\).  Eliminating \(x_1,x_2,x_3\) from the three first-derivative equations gives a homogeneous rational function of the coefficients, defined up to a numerical factor, and this function is called the discriminant of the form \(f\).  Its explicit calculation is possible only in the simplest cases.

If infinitely many singular points occur, then the partial derivatives have a common factor.  Conversely, a common factor of the first partial derivatives gives a component along which all these derivatives vanish.  Thus multiple components and higher singularities are detected algebraically by common divisors of the derivative forms.

Let \(P\) be a non-singular point of the curve.  Choose coordinates so that \(P\) is the vertex \(I_3\).  Then, expanding \(f\) according to the powers of \(x_3\), one has
\[
f=f_0x_3^n+f_1x_3^{n-1}+f_2x_3^{n-2}+\cdots ,
\]
where \(f_i\) is a binary form in \(x_1,x_2\).  Since \(P\) lies on the curve, \(f_0=0\), and since \(P\) is non-singular, \(f_1=0\) is the tangent at \(P\).  If this tangent is chosen to be \(x_1=0\), then the equation begins in the form
\[
f=x_1x_3^{n-1}+x_3^{n-2}f_2+\cdots .
\]
If the quadratic form \(f_2\) is divisible by \(x_1\), then the tangent \(x_1=0\) meets the curve at \(P\) with multiplicity at least three.  The point \(P\) is then called an inflection point, and the tangent is called the inflection tangent.

If \(P\) is an inflection point and \(x_1=0\) its inflection tangent, the form \(f\) may be written
\[
f=x_1\varphi+x_3^3\psi,
\]
where \(\varphi\) has degree \(n-1\) and \(\psi\) degree \(n-3\).  Conversely, any form of this shape has \(P\) as an inflection point and \(x_1=0\) as its inflection tangent.

A line touching the curve in two distinct points is called a bitangent.  If \(x_3=0\) is a bitangent and its two points of contact are the two vertices lying on this line, then, on putting \(x_3=0\), the remaining binary form must have \(x_1^2x_2^2\) as a factor.  Thus \(f\) has the form
\[
f=x_3\varphi+x_1^2x_2^2\psi,
\]
with \(\varphi\) of degree \(n-1\) and \(\psi\) of degree \(n-4\).  More generally, if \(x_3=0\) is a bitangent, the restriction of \(f\) to that line is the square of a quadratic factor; the two points of contact are the intersections of \(x_3=0\) with the corresponding conic.

\webersection{87}{Fundamental Covariants of a Ternary Form}

We recall the notion of invariants and covariants.  Suppose the ternary form \(f(x)\), with coefficients \(a\), is transformed by a linear substitution into \(\varphi(y)\), with coefficients \(b\), and let \(r\) be the determinant of the substitution.  A form
\[
\Phi(x,a)
\]
homogeneous in the variables and in the coefficients is a covariant if it satisfies a transformation law
\[
\Phi(y,b)=r^\lambda\Phi(x,a).
\]
The integer \(\lambda\) is called the weight.  If \(\Phi\) is independent of the variables, it is an invariant.

For every form of degree greater than one there is a covariant of weight \(2\), the Hessian determinant:
\[
H=\det\left(\frac{\partial^2 f}{\partial x_i\,\partial x_j}\right)_{1\le i,j\le 3}.
\]
It is of degree \(3(n-2)\) in the variables and of degree \(3\) in the coefficients.

Starting from the Hessian \(H\), one obtains further fundamental covariants by forming bordered determinants with \(f\), \(H\), and their first derivatives.  In Weber's notation these are denoted \(J\) and \(K\).  They are covariants of the original ternary form, of successively higher order in the variables and higher degree in the coefficients.  For a ternary form of degree \(n\) their orders and coefficient-degrees are, respectively,
\[
\begin{array}{c|c|c}
\text{covariant} & \text{order in }x & \text{degree in coefficients}\\
\hline
H & 3n-6 & 3\\
J & 6n-18 & 8\\
K & 12n-21 & 12.
\end{array}
\]
Only the first of these, the Hessian curve, is immediately needed for the elementary geometry of inflection points; the later ones enter in the reduction of the cubic form to its canonical form.

\webersection{88}{The Hessian Curve}

The covariant \(H\) has a simple geometrical meaning.  Since \(H\) is invariantly attached to \(f\), no generality is lost by choosing the coordinate triangle conveniently.  Put a vertex \(I_3\) at a non-singular point of the curve and choose the tangent there as the line \(x_1=0\).  Then the equation of the curve has the local form
\[
f=hx_1x_3^{n-1}
+x_3^{n-2}(ax_1^2+2bx_1x_2+cx_2^2)
+\cdots ,
\]
with \(h\ne0\).  Calculating the Hessian in descending powers of \(x_3\), the constant term at \(I_3\) is proportional to \(c\).  Consequently the Hessian curve passes through \(I_3\) precisely when \(c=0\), that is, precisely when \(I_3\) is an inflection point of \(f\).

Furthermore, the inflection tangent is also tangent to the Hessian exactly when the contact is of order four.  For irreducible curves of the third order this exceptional case cannot occur unless the curve has a factor; hence in the non-singular irreducible cubic case no such coalescence disturbs the count.

The Hessian curve also passes through the singular points of \(f\).  If we restrict attention to a curve without singular points, the intersection points of \(f=0\) with its Hessian \(H=0\) are exactly the inflection points of \(f\).  By Bézout's theorem a curve of order \(n\) without singular points has
\[
n\cdot 3(n-2)=3n(n-2)
\]
inflection points, counted with multiplicity.  In particular:

A non-singular curve of the third order has nine distinct inflection points.

\webersection{89}{Inflection Points of a Curve of the Third Order}

Let \(f=0\) now be a non-singular cubic curve.  Place two vertices of the coordinate triangle at two of its nine inflection points and choose the corresponding inflection tangents as coordinate sides.  Then, after a suitable normalization, the cubic may be brought to the form
\[
h(x_1^3+x_2^3+x_3^3)+6mx_1x_2x_3=0,
\]
or, after a linear change of notation,
\[
y_1^3+y_2^3+y_3^3+6my_1y_2y_3=0.
\]
This is the canonical form of the ternary cubic.  Each coordinate side contains three inflection points.  Conversely, the existence of one such triple of inflection points on a line allows the form to be arranged in this way.

For a non-singular cubic the nine inflection points have a highly regular configuration.  Given any two of them, there is a third such that the three lie on a line.  Thus the inflection points may be arranged into triples.  There are twelve such triples in all: every one of the nine points occurs in four triples, and counting incidences gives
\[
\frac{9\cdot4}{3}=12.
\]
This incidence structure will become the basis of the ``triple equations'' below.

In the canonical form the nine inflection points can be written explicitly.  If \(\varepsilon\) is a primitive cube root of unity, they are
\[
\begin{array}{ccc}
(0,1,-1),&(0,1,-\varepsilon),&(0,1,-\varepsilon^2),\\
(-1,0,1),&(-\varepsilon,0,1),&(-\varepsilon^2,0,1),\\
(1,-1,0),&(1,-\varepsilon,0),&(1,-\varepsilon^2,0).
\end{array}
\]
Thus three of them are real, while the other six are imaginary, in conjugate pairs, when the ground field is real and the canonical parameter is real.

The twelve lines through triples can be divided in four ways into three lines so that each of the nine inflection points occurs exactly once in each division.  A convenient notation is to label the nine points by pairs
\[
(\xi,\eta),\qquad \xi,\eta\in\mathbb Z/3\mathbb Z,
\]
with conjugation interchanging \((\xi,\eta)\) and \((\eta,\xi)\).  Then three labelled points lie on one of the twelve lines exactly when
\[
\xi_1+\xi_2+\xi_3\equiv0,\qquad
\eta_1+\eta_2+\eta_3\equiv0\pmod 3.
\]
This arithmetical notation makes the incidence system transparent.

\webersection{90}{Transformation of the Cubic Form to the Canonical Form}

To reduce a ternary cubic to the canonical form, one must express the required substitution algebraically.  For this purpose Weber computes the fundamental covariants of the canonical cubic
\[
\varphi(y)=h(y_1^3+y_2^3+y_3^3)+6my_1y_2y_3.
\]
The Hessian covariant, after multiplication by a convenient numerical factor, is again a cubic of the same shape:
\[
\Delta(y)=-hm^2(y_1^3+y_2^3+y_3^3)+(h^3+2m^3)y_1y_2y_3.
\]
A second covariant \(J\), formed from a bordered determinant involving \(f\), \(\Delta\), and their derivatives, is likewise evaluated explicitly.  From \(f,\Delta,J\) one can express the elementary symmetric functions
\[
y_1y_2y_3,\qquad
y_1^3+y_2^3+y_3^3,\qquad
y_1^3y_2^3+y_2^3y_3^3+y_3^3y_1^3
\]
in terms of the corresponding covariants of the original cubic.

The decisive equations may be summarized as follows.  If \(r\) is the determinant of the transformation and the coefficient of \(y_1^3+y_2^3+y_3^3\) has been normalized to \(h=1\), then
\[
(1+8m^3)y_1y_2y_3=m^2 f+r^2\Delta,
\]
and a second formula gives \(y_1^3+y_2^3+y_3^3\) in terms of \(f\) and \(\Delta\).  A third covariant \(K\) gives
\[
r^9K=(1+8m^3)^2
(y_1^3-y_2^3)(y_2^3-y_3^3)(y_3^3-y_1^3).
\]
The factor \(1+8m^3\) cannot vanish when the cubic is non-singular; indeed in that case the three first partial derivatives would have a common zero.

Thus, once \(m\) and \(r\) have been determined, the quantities \(y_1^3,y_2^3,y_3^3\) are the roots of a cubic equation
\[
u^3-Qu^2+Pu-R^3=0,
\]
whose coefficients are rational functions of \(f,\Delta,J\) and of the invariants \(m,r\).  The extraction of cube roots from its roots then gives the linear forms \(y_1,y_2,y_3\), and hence the desired substitution.  The reduction of the original cubic to canonical form is therefore converted into the solution of certain algebraic equations attached to its covariants and invariants.

\webersection{91}{The Invariants of the Curve of the Third Order and the Biquadratic Equation}

It remains to find the invariants on which the canonical parameter \(m\) depends.  Consider the pencil of cubic covariants
\[
\Delta_{\lambda,\mu}=\lambda f+\mu\Delta.
\]
The Hessian of this pencil can again be expressed as a linear combination of \(f\) and \(\Delta\):
\[
H(\Delta_{\lambda,\mu})=L(\lambda,\mu)f+M(\lambda,\mu)\Delta.
\]
The coefficients of the binary forms \(L\) and \(M\) are invariants of the original cubic.  They may be found by comparing the coefficients of the ten monomials in a ternary cubic.  If
\[
a_i,\quad A_i,\quad B_i
\]
are the corresponding coefficients of \(f,\Delta,H(\Delta_{\lambda,\mu})\), then
\[
La_i+MA_i=B_i,
\]
and hence, for distinct indices \(i,k\),
\[
L(a_iA_k-a_kA_i)=B_iA_k-B_kA_i,\qquad
M(a_iA_k-a_kA_i)=-B_ia_k+B_ka_i.
\]
This also shows that \(L\) and \(M\) are integral functions of the coefficients of \(f\); a denominator would have to divide all determinants \(a_iA_k-a_kA_i\), which is impossible, as is seen from a special cubic.

For the canonical form the calculation gives two fundamental invariants, of degrees \(4\) and \(6\), denoted by \(S\) and \(T\).  With \(h=1\) they satisfy
\[
m(1-m^3)=r^4S,
\]
and
\[
1-20m^3-8m^6=r^6T.
\]
Eliminating the scale factor \(r\) yields a biquadratic equation for the appropriate expression in \(m\).  This equation is the algebraic obstacle in the reduction of the cubic to the canonical form; after adjoining one of its roots the remaining steps consist of the cubic equation for \(y_1^3,y_2^3,y_3^3\) and the extraction of cube roots.

The inflection points are also visible at this stage.  In the canonical coordinates they are obtained by putting successively one of \(y_1,y_2,y_3\) equal to zero in
\[
y_1^3+y_2^3+y_3^3+6my_1y_2y_3=0.
\]
With \(\varepsilon^3=1\), \(\varepsilon\ne1\), the nine points are precisely the nine triples displayed above.  Their twelve lines are governed by the congruence rule
\[
\xi_1+\xi_2+\xi_3\equiv
\eta_1+\eta_2+\eta_3\equiv0\pmod3.
\]
The real form of this configuration is as follows.  The three points
\[
(0,0),\quad(1,1),\quad(2,2)
\]
are real, and the remaining six are divided into three conjugate pairs.  Among the twelve lines, some are real and some occur in conjugate imaginary pairs, exactly as dictated by the table of triples.

\webersection{92}{Triple Equations}

The nine inflection points of a non-singular cubic give rise to a purely algebraic structure.  Let
\[
x_1,\dots,x_9
\]
be nine quantities arranged in twelve triples, with the property that any two quantities determine exactly one triple and every quantity lies in four triples.  If, whenever \(x_1,x_2,x_3\) are the members of a triple, there is a rational relation
\[
x_1=\Phi(x_2,x_3),
\]
or equivalently each member of a triple is rationally determined by the other two, then an equation whose nine roots are so arranged is called a triple equation.

The incidence system of the nine inflection points shows that such a triple equation is naturally labelled by pairs
\[
(\xi,\eta)\in(\mathbb Z/3\mathbb Z)^2.
\]
The triples are exactly the sets
\[
(\xi_1,\eta_1),\quad(\xi_2,\eta_2),\quad(\xi_3,\eta_3)
\]
for which
\[
\xi_1+\xi_2+\xi_3\equiv0,\qquad
\eta_1+\eta_2+\eta_3\equiv0\pmod3.
\]
Conversely, beginning with any one triple and one fourth point, this rule reconstructs the entire system of twelve triples.  Thus the geometry of the flexes of a cubic is transferred to the algebra of an equation of degree \(9\).

\webersection{93}{The Group of the Triple Equations}

We now determine the Galois group \(P\) of a triple equation.  A permutation in \(P\) must carry every triple into a triple.  Passing to the notation
\[
x=(\xi,\eta),
\]
this means that a permutation must preserve the rule
\[
\xi_1+\xi_2+\xi_3\equiv0,\qquad
\eta_1+\eta_2+\eta_3\equiv0\pmod3.
\]

Every permutation of the nine pairs can be represented by functions
\[
\xi'=\phi(\xi,\eta),\qquad \eta'=\psi(\xi,\eta)\pmod3,
\]
where the degree in each variable is at most \(2\).  Holding \(\xi\) fixed and applying the permutation to the triple
\[
(\xi,0),\quad(\xi,1),\quad(\xi,2),
\]
one sees that the quadratic terms in \(\eta\) must vanish.  The same argument with the variables interchanged removes the quadratic terms in \(\xi\).  Applying the condition to the diagonal triple
\[
(0,0),\quad(1,1),\quad(2,2)
\]
then removes the mixed quadratic term.  Consequently every substitution in \(P\) has the affine linear form
\[
\begin{aligned}
\xi'&=a\xi+b\eta+c,\\
\eta'&=a'\xi+b'\eta+c'
\end{aligned}
\qquad(\bmod 3),
\]
with non-zero determinant
\[
ab'-a'b\not\equiv0\pmod3.
\]
Thus the group of a triple equation is contained in the full affine linear congruence group over \(\mathbb F_3\).

Conversely, every substitution of this affine linear form carries triples into triples, since it preserves the vanishing of the sums of the two coordinates.  Hence a degree-nine equation whose Galois group is a doubly transitive subgroup of this affine group is a triple equation.  The affine group decomposes as
\[
P=SQ,
\]
where \(Q\) is the translation group
\[
\xi'=\xi+\alpha,\qquad \eta'=\eta+\beta
\]
and \(S=\GL(2,3)\) is the homogeneous part.

The translation group \(Q\) is simply transitive: if one root is fixed by a translation, all are fixed.  The stabilizer of \((0,0)\) in \(P\) is the intersection of \(P\) with \(S\).  Using the composition series of \(S\) discussed earlier, Weber shows that the smaller subgroups \(S_4Q\) and its immediate enlargements are still not doubly transitive; true triple equations begin only at the subgroup \(S_2Q\), and therefore also at \(S_1Q\) and \(SQ\).  Certain other affine subgroups, though transitive, are not groups of genuine triple equations.

\webersection{94}{Reality Relations of the Triple Equations}

Assume now that the rationality field is real.  If a triple contains two real roots, the relation defining the third root shows that the third is also real.  If an imaginary root occurs, then in every one of the four triples containing it there must be at least one further imaginary root.  Hence the number of imaginary roots is at least \(5\), and therefore, since imaginary roots occur in conjugate pairs, at least \(6\).

A triple containing two real roots, or two conjugate imaginary roots, will be called a real triple.  If a triple contains imaginary roots, then the conjugate roots form the conjugate triple.

There are therefore only three possible reality patterns:
\[
\begin{array}{ll}
1.&\text{all nine roots real},\\
2.&\text{one real root and eight imaginary roots},\\
3.&\text{three real roots and six imaginary roots}.
\end{array}
\]
In the second case the four conjugate pairs must form four triples, all having the same real third root.  If this real root is labelled \((0,0)\), then the conjugate pairs are
\[
(1,2)(2,1),\quad
(1,1)(2,2),\quad
(1,0)(2,0),\quad
(0,1)(0,2).
\]

In the third case the three real roots themselves form a triple.  We may take them to be
\[
(0,0),\quad(1,1),\quad(2,2).
\]
The conjugate imaginary pairs cannot lie in the same row of the remaining arrangement, since otherwise the third root of that row would be real.  A short check of the twelve triples shows that each of the three real roots occurs, besides the real triple just displayed, in one and only one further real triple whose other two roots are conjugate imaginary.  This gives exactly the same arrangement of real and imaginary points as in the inflection points of a real cubic curve.

The first and second cases also occur.  Indeed, from a general equation of the ninth degree one obtains a triple equation by adjoining a function belonging to the affine linear group \(SQ\).  Weber exhibits such a function as the sum of the products of the three roots over the twelve triples.  This function is real in all three reality cases, and hence the enlarged rationality field in which the equation becomes a triple equation is itself real.




\booktitle{Twelfth Section. The Bitangents of a Curve of the Fourth Order}

\webersection{95}{Number of Bitangents of a Curve of the Fourth Order}

A geometrical problem of particular importance, both because of its relations to other fields and because it presents algebraic phenomena analogous to those arising from the inflection points of cubic curves, is the problem of the bitangents of a plane quartic.  A bitangent is a line touching the curve at two distinct points.  For curves of order less than four there are no bitangents; for a quartic a bitangent has no further intersection with the curve.  In special cases the two points of contact may coincide, producing a tangent of fourfold contact.

Let
\[
f(x_1,x_2,x_3)=0
\]
be the equation of a non-singular curve of the fourth order.  If \(x\) and \(y\) are two fixed points, the points on the line joining them have coordinates \(x+ty\).  Substituting gives
\[
f(x+ty)=0,
\]
a quartic equation in \(t\).  Written according to powers of \(t\), it has the form
\[
f(x+ty)=P_0+4tP_1+6t^2P_2+4t^3P_3+t^4P_4,
\]
where \(P_\nu(x,y)\) are the polars of \(f\).

If \(x\) lies on the curve, \(P_0=0\).  If \(y\) lies on the tangent at \(x\), then \(P_1=0\).  The two remaining intersections of the tangent line with the quartic are determined by
\[
6P_2+4tP_3+t^2P_4=0.
\]
The tangent is a bitangent precisely when this quadratic has a double root.  Its discriminant must vanish; in Weber's notation the condition is
\[
R(x,y)=2P_3^2-3P_2P_4=0.
\]
Thus \(R(x,y)=0\), with \(x\) on the curve and \(y\) on the tangent at \(x\), is the necessary and sufficient condition for \(x\) to be a point of contact of a bitangent.

Weber then eliminates \(y\) from this condition and the tangent equation.  The result is a curve
\[
\chi(x)=0
\]
of degree \(14\) which passes through the contact points of the bitangents and through no other point of the quartic.  Since the quartic \(f=0\) and the curve \(\chi=0\) meet in
\[
4\cdot14=56
\]
points, and since every bitangent contributes two contact points, the number of bitangents is
\[
56/2=28.
\]
Possible exceptional contacts are excluded by specializing to a quartic for which the contacts are distinct and then using the fact that the number cannot change while the curve remains non-singular.

Therefore a non-singular plane quartic has twenty-eight bitangents.

\webersection{96}{The Steiner Complexes}

Take two bitangents \(x_1=0\), \(y_1=0\).  Since their four points of contact lie on a conic, the equation of the quartic may be written in the form
\[
f=x_1y_1v-u^2,
\]
where \(u=0\) is a conic and \(v=0\) another conic.  If a second representation
\[
f=x_2y_2v_1-u_1^2
\]
is possible, then comparison gives
\[
x_1y_1(v-v_1)=(u-u_1)(u+u_1).
\]
One of the two factors \(u-u_1\), \(u+u_1\) must be divisible by \(x_1y_1\).  After changing the sign of \(u_1\) if necessary,
\[
u_1=u+\lambda x_1y_1,
\]
and hence
\[
f=x_1y_1\bigl(v+2\lambda u+\lambda^2x_1y_1\bigr)
  -(u+\lambda x_1y_1)^2.
\]
Whenever the quadratic form
\[
v+2\lambda u+\lambda^2x_1y_1
\]
splits into two linear factors \(x_2y_2\), the two lines \(x_2=0\), \(y_2=0\) are two further bitangents, and the eight points of contact of
\[
x_1,\ y_1,\ x_2,\ y_2
\]
lie on a conic.

The condition that this conic split is the vanishing of its determinant.  It is an equation of the fifth degree in \(\lambda\), and therefore yields five further pairs of bitangents.  Hence:

To every pair of bitangents \(x_1,y_1\) belong five further pairs \(x_i,y_i\) such that the eight contact points of \(x_1,y_1,x_i,y_i\) lie on a conic.

The six pairs
\[
x_1y_1,\ x_2y_2,\ x_3y_3,\ x_4y_4,\ x_5y_5,\ x_6y_6
\]
so determined form what Weber calls a Steiner complex.

Two bitangents in a complex form a syzygetic pair; three bitangents are called syzygetic if their six contact points lie on a conic, and azygetic otherwise.  Weber proves that three bitangents of one Steiner complex, no two of which form one of the distinguished pairs \(x_i,y_i\), are always azygetic.  This follows directly from the equation
\[
f=x_1y_1x_2y_2-u_1^2
\]
and from the impossibility, for a non-singular quartic, that the relevant three lines meet in a point or that a conic pass through the wrong six contact points.

\webersection{97}{Pairs and Triples of Complexes}

The preceding results provide an effective calculus for the twenty-eight bitangents.  Begin with one Steiner complex
\[
x_1y_1,\ x_2y_2,\ \dots,\ x_6y_6.
\]
The form of the equation shows that the complex determined by the pair \(x_1,x_2\) contains the pair \(y_1,y_2\), while the complex determined by \(x_1,y_2\) contains \(x_2,y_1\).  These two new complexes, together with the original one, contain all twenty-eight bitangents.

Consequently every syzygetic triple of bitangents is completed by one uniquely determined fourth bitangent to a syzygetic quadruple; equivalently, a conic through the six contact points of three such bitangents cuts the quartic again in the two contact points of a fourth bitangent.

Counting pairs gives the number of complexes.  There are
\[
\binom{28}{2}=14\cdot27
\]
pairs of bitangents.  Each Steiner complex contains six distinguished pairs.  Since every pair determines one complex, the total number of Steiner complexes is
\[
\frac{14\cdot27}{6}=63.
\]

Two complexes are called syzygetic if they have four pairs of bitangents in common, and azygetic if they have six azygetic elements in common.  A syzygetic pair of complexes is completed in one and only one way to a syzygetic triple of complexes; an azygetic pair is similarly completed to an azygetic triple.  In a syzygetic triple all twenty-eight bitangents occur, while in an azygetic triple only eighteen occur.

\webersection{98}{The Aronhold Seven-Systems}

Especially important are azygetic systems of seven bitangents, first considered by Aronhold.  Weber calls them Aronhold seven-systems, complete seven-systems, or simply complete systems.

Such systems exist.  For instance, from an azygetic pair of complexes one obtains a set of seven bitangents no three of which form a syzygetic triple.  Conversely, no azygetic system can contain more than seven elements.

The key theorem is:

Any six elements of a complete seven-system occur in one and only one Steiner complex.

Uniqueness is proved by contradiction: if the six elements belonged to two different complexes, the corresponding syzygetic triple of complexes would force two of the six to be paired inside a third complex, which would create a forbidden syzygetic triple.  Existence is obtained by choosing two of the six, forming a complex through them, and using the decomposition of the associated syzygetic complex triple to show that the other four must enter the same complex.

This theorem makes it possible to introduce a notation for all twenty-eight bitangents.  Let a complete system be
\[
x_1,x_2,\ldots,x_7.
\]
Introduce an eighth symbol \(8\), and denote the seven elements by
\[
[18],[28],\ldots,[78].
\]
Then all twenty-eight bitangents are denoted uniformly by
\[
[\mu\nu],\qquad 1\le \mu<\nu\le8.
\]
This is Cayley's notation.

\webersection{99}{The Hesse--Cayley Algorithm for Denoting the Bitangents}

The notation \([\mu\nu]\) is best visualized as a line segment joining two of eight labelled points.  Two symbols with no common index, such as \([12]\) and \([34]\), are represented by disjoint segments; two symbols with one common index, such as \([12]\) and \([13]\), by two segments issuing from one point.

In this notation the distinction between syzygetic and azygetic triples becomes combinatorial.  Weber checks the possible types of three segments.  Triples shaped like a star, a path, or a triangle of certain types are azygetic; the remaining types are syzygetic.  The check is made by comparing them with the three complexes obtained from an azygetic complex triple.

The complete seven-systems have only two apparent graphical shapes:
\[
\text{a seven-rayed star},\qquad
\text{a triangle together with a four-rayed star}.
\]
There are eight systems of the first type, one for each possible centre.  For the second type, the triangle can be chosen in
\[
\frac{8\cdot7\cdot6}{2\cdot3}=56
\]
ways, and then the centre of the four-rayed star in five ways, giving \(280\).  Thus the total number of Aronhold systems is
\[
8+280=288.
\]
The two graphical shapes do not represent two intrinsically different species; the difference lies only in the notation.

In the same notation the Steiner complexes have two standard forms.  One is
\[
[12][34],\ [13][24],\ [14][23],\
[56][78],\ [57][68],\ [58][67],
\]
and the other is
\[
[17][18],\ [27][28],\ [37][38],\
[47][48],\ [57][58],\ [67][68],
\]
with the obvious relabellings.

\webersection{100}{Rational Determination of the Curve from a Complete Seven-System}

The importance of Aronhold systems is expressed in two facts.

First, if a complete seven-system of bitangents of a quartic is known, then all the remaining bitangents can be determined rationally from it.

Second, given seven general lines in the plane, one can, provided certain rational functions of their coefficients do not vanish, determine rationally the equation of a non-singular quartic for which these seven lines form an Aronhold system.

To prove the first assertion, it suffices by symmetry to determine three new bitangents at once.  Let
\[
x_1,\ldots,x_7
\]
be the given complete system.  Form the three Steiner complexes in which six of the seven occur, excluding in turn \(x_1,x_2,x_3\).  The newly appearing bitangents are denoted \(\ell_1,\ell_2,\ell_3,\ldots\).  By choosing suitable constant factors one may write the quartic in the form
\[
\sqrt{x_1\ell_1}+\sqrt{x_2\ell_2}+\sqrt{x_3\ell_3}=0,
\]
or, equivalently,
\[
f=4x_2\ell_2x_3\ell_3-u_1^2
 =4x_3\ell_3x_1\ell_1-u_2^2
 =4x_1\ell_1x_2\ell_2-u_3^2,
\]
with
\[
u_1=-x_1\ell_1+x_2\ell_2+x_3\ell_3
\]
and the two analogous formulae obtained cyclically.  The syzygetic relations among the complexes yield linear equations from which \(\ell_1,\ell_2,\ell_3\) are rationally determined.  Repeating the same argument determines all remaining bitangents.

The second assertion is obtained by reversing the construction.  From seven general lines one forms the rational functions which play the role of the missing \(\ell\)'s, and then the displayed equation gives the quartic.  The exceptional cases occur exactly when the determinants used in solving the linear systems vanish, or when the resulting quartic becomes singular.

\webersection{101}{The Galois Group of the Bitangent Problem}

The bitangents are roots of an equation of degree \(28\).  Choose a line \(L\) not tangent to the quartic and not passing through any intersection point of two bitangents.  In Cartesian coordinates with \(L\) as the \(x\)-axis, a bitangent can be written in the form
\[
y=\theta(x-\xi).
\]
If \(F(x,y)=0\) is the equation of the quartic, the condition that
\[
F(x,\theta(x-\xi))
\]
be a perfect square quartic in \(x\) gives two equations in \(\xi,\theta\).  Eliminating \(\theta\) gives the bitangent equation of degree \(28\).

The roots of this equation inherit all the arrangements of the corresponding bitangents: Steiner complexes, Aronhold systems, syzygetic and azygetic triples.  Denote the roots, as before, by symbols \([ik]\), \(1\le i<k\le8\).

Given two roots \(p,q\), one can rationally construct the conic through their four contact points and therefore form
\[
f=pqv-u^2.
\]
The five conics in the pencil
\[
v+2\lambda u+\lambda^2pq=0
\]
which split into pairs of lines give the ten other bitangents in the Steiner complex determined by \(p,q\).  Hence there is a rational function
\[
X(\xi_1,\xi_2,\xi_3)
\]
which vanishes exactly when \(\xi_1,\xi_2,\xi_3\) form a syzygetic triple, and does not vanish for an azygetic triple.

Similarly, by adjoining a complete seven-system
\[
\xi_1,\ldots,\xi_7
\]
all other roots are rationally determined.  Therefore every permutation in the Galois group must preserve the whole incidence structure of syzygetic triples and complete seven-systems.  Conversely, every permutation preserving this structure is realized in the group.  Thus the Galois group of the bitangent equation is the full group of permutations of the twenty-eight symbols \([ik]\) induced by the transformations described in the next section.

Counting the complete systems gives the order of the group.  There are \(288\) choices of a complete seven-system, and each such system can be arranged in \(7!\) ways.  Thus
\[
|P|=288\cdot7!
   =8!\cdot36.
\]

\webersection{102}{Representation of the Group}

The notation \([ik]\) immediately shows that every permutation of the eight indices \(1,\dots,8\) gives a permutation of the twenty-eight roots.  These permutations form a subgroup
\[
S\cong S_8
\]
of index \(36\) in the full group \(P\).

The missing cosets are obtained by replacing the original Aronhold system
\[
[18],[28],\ldots,[78]
\]
by a system of the other type, for instance
\[
[23],[31],[12],[48],[58],[68],[78].
\]
The induced permutation is denoted
\[
\Pi_{1,2,3,8}=\Pi_{4,5,6,7},
\]
where the subscript may be replaced by any set of four indices or by its complement.  There are \(35\) such permutations.

The whole group is therefore
\[
P=S+\sum S\Pi_{\alpha_1\alpha_2\alpha_3\alpha_4},
\]
where the sum extends over the \(35\) partitions of the eight indices into two complementary sets of four.

The action of \(\Pi_{\alpha_1\alpha_2\alpha_3\alpha_4}\) is simple.  If the two indices of \([\mu\nu]\) both lie in the same half of the partition, \([\mu\nu]\) is changed to the complementary pair in that half; if one index lies in each half, \([\mu\nu]\) remains fixed.  Each \(\Pi\) has order two.

The composition laws are consequently explicit.  If \(\sigma\in S_8\), then
\[
\Pi_{1234}\sigma=\sigma\Pi_{\sigma(1)\sigma(2)\sigma(3)\sigma(4)},
\]
and
\[
\Pi_{1234}^2=1.
\]
For overlapping quadruples one obtains formulae such as
\[
\Pi_{1234}\Pi_{1235}=(45)\Pi_{1234},
\]
and
\[
\Pi_{1234}\Pi_{1256}
=(12)(34)(56)(78)\Pi_{1278}.
\]
These relations completely determine the group \(P\).

\webersection{103}{Simplicity of the Group of the Bitangent Problem}

The representation just obtained gives a direct proof that \(P\) is simple.  Write
\[
P=S+\sum S\Pi_{\alpha_1\alpha_2\alpha_3\alpha_4},
\]
where \(S\cong S_8\).  Let \(S'\) be the alternating group on the eight letters.  Every non-trivial normal subgroup of \(S_8\) contains \(S'\).

Let \(Q\) be a non-trivial normal divisor of \(P\).  If \(Q\cap S\) contains a non-identity element, then it contains \(S'\).  But if \(Q\) contains \(S'\), conjugating by the \(\Pi\)'s and using the relations of §102 gives elements in all cosets \(S\Pi_{\alpha_1\alpha_2\alpha_3\alpha_4}\).  Thus \(Q=P\).

It remains to exclude the case in which \(Q\) contains no non-trivial element of \(S\).  If \(Q\) contains an element \(\sigma\Pi_{1234}\), conjugate it by a transposition \((\alpha\beta)\) chosen so that the conjugated transposition differs from the original one.  Multiplying the two conjugates produces a non-identity element of \(S\), contradicting the assumption unless \(Q=P\).  Therefore the only normal divisors of \(P\) are \(1\) and \(P\).

Consequently the group of the bitangent equation is simple.

An immediate consequence is that the group acts by even permutations on the twenty-eight roots.  Hence the discriminant of the bitangent equation is a square.  The stabilizer of one root is transitive on the other \(27\), so \(P\) is doubly transitive, but it is not triply transitive: after two roots are fixed, a root syzygetic with them cannot be carried to an azygetic one.  Adjoining two roots splits off a factor of degree \(10\), whose roots are the bitangents completing the Steiner complex with the two given roots.

The subgroup \(S\cong S_8\), of index \(36\), is also distinguished.  By adjoining a root of a degree-\(36\) equation one reduces the bitangent problem to a general equation of degree \(8\).  In this sense the bitangent problem is governed by an \(8\)-letter symmetric group sitting inside the simple group \(P\) of order \(8!\cdot36\).




\webersection{104}{Reality of the Bitangents}

Let the quartic curve be real and non-singular.  If a bitangent \(l\) is imaginary, its conjugate \(l'\) is also a bitangent.  Every rational relation among the roots of the bitangent equation remains valid after replacing all imaginary roots by their conjugates.  Hence conjugation carries syzygetic systems to syzygetic systems and azygetic systems to azygetic systems.

A Steiner complex \(C\) is therefore carried by conjugation into a Steiner complex \(C'\).  If \(C=C'\), the complex is called real; otherwise \(C,C'\) form a conjugate pair of complexes.  In a real complex a pair is either a pair of real bitangents, a conjugate pair \(l,l'\), or part of a conjugate double-pair
\[
(l,m),\quad(l',m').
\]
A real bitangent can never be paired in a real complex with a non-real bitangent, since then the conjugate pair would share the same real element, which is impossible inside a Steiner complex.

Similarly, an Aronhold seven-system is carried by conjugation into another Aronhold system.  It is real if it is identical with its conjugate; such a system contains, together with every imaginary bitangent, its conjugate, and therefore contains an odd number of real bitangents.

Weber derives the following restrictions.

First, not all twenty-eight bitangents can be imaginary.  Starting with a conjugate pair and a third imaginary bitangent, one would obtain conjugate Steiner complexes arranged in a way that forces the existence of a real pair.  Thus at least one real pair occurs, and then further incidence arguments force at least four real bitangents.

Second, there are always four real bitangents forming two pairs in a real Steiner complex.

Third, if one further real bitangent exists beyond these four, then there are eight real bitangents, grouped as four pairs in one Steiner complex.

Fourth, if more than eight real bitangents occur, then there are sixteen real bitangents arranged in a syzygetic triple of Steiner complexes, each complex containing four real pairs.

Repeating the argument shows that the only possible numbers and arrangements of real bitangents of a real non-singular quartic are:
\[
\begin{array}{ll}
1.&4\text{ real syzygetic bitangents},\\
2.&8\text{ real bitangents, forming four pairs of a Steiner complex},\\
3.&16\text{ real bitangents, arranged in a syzygetic complex triple},\\
4.&28\text{ real bitangents}.
\end{array}
\]

\webersection{105}{Proof of the Existence of the Four Cases}

The previous section proved only that no other cases are possible.  To show that all four really occur, Weber uses the theorem of §100: from seven general lines one can construct rationally a non-singular quartic for which those seven lines are an Aronhold system, and from this system all bitangents are rationally determined.

The seven given lines may include imaginary conjugate pairs.  The resulting curve is real whenever the conjugate of the seven-system is again an Aronhold system of the same curve.  This is certainly the case when the seven-system itself is real, i.e. when every imaginary line occurs together with its conjugate.

A real Aronhold system of seven lines can contain:
\[
\begin{array}{ll}
1.&1\text{ real line and }3\text{ conjugate imaginary pairs},\\
2.&3\text{ real lines and }2\text{ conjugate imaginary pairs},\\
3.&5\text{ real lines and }1\text{ conjugate imaginary pair},\\
4.&7\text{ real lines}.
\end{array}
\]
Using Cayley's notation for the bitangents, the remaining twenty-one bitangents are computed from the seven by the Hesse--Cayley rules.  In the first case, if the given lines are
\[
0,\ 1,1',\ 2,2',\ 3,3',
\]
then the real bitangents are exactly
\[
0,\ [11'],\ [22'],\ [33'],
\]
and all others are imaginary.  This gives the case of four real bitangents.

In the second case, with
\[
0,\ 1,\ 2,\ 3,3',\ 4,4',
\]
the real bitangents are
\[
0,\ 1,\ 2,\ [12],\ [20],\ [01],\ [33'],\ [44'],
\]
and the remaining bitangents occur in conjugate pairs.  These eight real bitangents form four pairs of a Steiner complex.

In the third case, with
\[
0,\ 1,\ 2,\ 3,\ 4,\ 5,5',
\]
the construction gives sixteen real bitangents.  They are arranged in a syzygetic triple of complexes, each containing four real pairs, and the remaining twelve are imaginary in conjugate pairs.

In the fourth case, when all seven lines of the Aronhold system are real, every bitangent determined rationally from them is real, and all twenty-eight bitangents are real.

This proves the existence of the four possible real cases.  Weber adds that the reality of a bitangent does not imply the reality of its two points of contact, since those contact points depend on a quadratic equation.  Thus the real bitangents themselves split further into bitangents with real contact points and bitangents with imaginary contact points; this finer classification is not pursued here.

\end{document}
